% !TeX root = ../papers/nc-p3.tex
\section{二次錐の全射影における標準辺の有限性と折り返しの排除}
\label{sec:cone-fold-plucker}
\begin{theorem}\label{thm:cone-fold-plucker}
二十次の整な一般被約余因子曲線を持つ被約窓の二次錐正規形を考える。
全ての特殊値で、標準 $m,y$ 辺は両射影で有限被約で、全二十印を与える。
両側の中間位置をともにこの $m$ にした四辺有限標準正則延長はない。
\end{theorem}
\begin{proof}[射影中心と平面切断の二項四次式]
全母線で入力は完全二次系を与える。そうでなければ豊富な入力束による有限射は
その母線を直線に二対一に写し、正規化押出し表示からその直線が階数二以下の軌跡に入る。
頂点の入力繊維は長さ二以下である。実際 $V$ の基底を $w^2$ の係数列が $(1,0,0,0)$
となるよう取り、残る三座標を $A_i(s,t)+w B_i(s,t)$ と書く。
母線の全階数により各 $z$ で三成分の $A(z),B(z)$ は独立である。ゆえに他の点は頂点の入力値に来ない。
また $B_i$ は基点のない二項二次系で階数二以上。階数三なら頂点の繊維長は一、
二なら残る接直線は接錐の二次式の零でない方向である。そうでなければ $B_i$ の共通根になる。
局所環の随伴次数環でこの二一次式を割ると残る変数の平方が零となり、長さは二以下である。
従って二十印では頂点を避けた長さ三繊維から定理\ref{thm:general-quadric-cone-exclusion}と同じ $m$ を作れる。

辺 $a m+bDm=0$ の任意の点では $b\ne0$。従って $y$ は頂点を通らない平面で、
$w=-a/b$ を代入すると $R/(y)=k[s,t]$ である。$V$ の二項四次式への制限像を $W$ とし、
これは基点を持たない。$m$ の解空間は $(S_5/W S_1)^*$ である。
$\dim W\le2$ ならこの平面切断円錐曲線が入力直線以下へ写り、有限性と次数四から
その直線が階数二曲線に入るので不可能。$\dim W=4$ なら $WS_1=S_5$。
実際零化五次汎関数の二収縮が同じ一直線に入れば、その係数のHankel二行の階数が一となる。
そのような汎関数は一点評価で、$W$ に基点ができる。
$\dim W=3$ なら $W\otimes S_1\to S_5$ は正方であり、余核があれば一次係数関係がある。
係数は二独立（さもなくば基底の定数関係）なので、基底変更で
$W=\langle s h,t h,g\rangle$、$h$ は三次、$g$ は $h$ と互いに素である。
従って $WS_1=hS_2+gS_1$ は五次元で、$y$ は実際に長さ三の同一入力繊維を含む。
その長さはちょうど三である。入力の長さは根行列の余階数に等しく、十成分が全二次式を張るので
階数零点はない。この繊維が張る平面は一意である。
階数一点は有限個（曲線 $C$ 上で一般階数は二）だから $y$ は有限個で各 $m$ は一意、
逆向きも $Dm/m$ により一意である。六双一次完全交叉の長さは二十で、既に二十の異なる印を含むから被約である。
\end{proof}

\begin{proof}[六つの八次Plücker式]
二行 $m(z),(Dm)(z)$ は全ての $z$ で独立であり、入力 $m$ 空間内の直線族 $\ell(z)$ を与える。
実際 $z$ で消える一次式を $l$ とし、対応母線上で消える非零な $v_z\in V$ を取る。
$h_z=v_z/l\in R_3$ は $l$ 倍で $R_4/V$ において零となる。評価列の双対同定から
$h_z$ は $S_5$ 上の $z$ での評価汎関数に比例する
（$l:S_4\to S_5$ の双対の核直線）。$z$ で消えない一次式 $u$ を取り、
$u^2 h_z$ は同様に $A$ で $S_3$ の評価を与える。
二行が従属なら $(w-c u^2)h_z\in V S_1$ となる。$l$ を掛けると
$V\otimes R_2\to R_6$ の非零核元を得る。$w$ の係数 $v_z$ は他の二項二次係数で消せない。
これは定理\ref{thm:quadric-kernel-projection}の積同型に反する。

折り返しの正則延長があると仮定する。有限辺の反転により頂テンソルは $m$ の二位置で対称である。
対称な三次交叉は線形写像 $L:S_3\to S_5^*$ で
$L(p)(a q)=L(q)(a p)$ および $L(p)(Dq)=L(q)(Dp)$ を満たすものになる。
実際 $R_{my}$ の六列は $S_5^*$ で、第一文字と最後の文字を入れ替えたときの係数を比較すればよい。
最初の等式は $L(p)(f)=\lambda(p f)$ と同値で、$\lambda\in S_8^*$ は一意である。
単項式の指数で $L_{i,j+k}=L_{j,i+k}$ ($0\le i,j\le3,0\le k\le2$) を課せば
二十四係数は和 $i+j+k$ ごとの九係数となる。ゆえに交叉は
\[
       W_D^\perp,\qquad W_D=\langle pDq-qDp:p,q\in S_3\rangle\subset S_8
\]
である。正則候補では四次元なので $\dim W_D=5$。この五空間は直線族 $\ell$ の六Plücker式の包でもある。
$\ell:\mathbb P^1\to\mathbb P(W_D^*)$ は基点を持たず、像は直線Grassmann多様体に入る
（消去二平面と直線の座標の双対同定を使う）。

対称反対切断の核に階数二以下の元があれば、二つの非零三次式の対称積 $p\odot q$ と書ける。
全 $\lambda\in W_D^\perp$ と全 $a\in S_2$ に対して $\lambda(a p q)=0$ なので
$pqS_2\subset W_D$。$\mathbb P(W_D^*)$ は四次元であり、この三条件は直線 $\Lambda$ を切る。
$\ell$ は長さ六の因子 $V(pq)$ をスキームとして $\Lambda$ に写すことになる。これは不可能である。
固定した母線像の逆像長は二以下である。その直線に乗る二つの独立な三次式は、
逆像スキーム上でともに $z$ における値が零になり、共通零の長さは二以下だからである。
従って $\Lambda$ がGrassmann二次超曲面に含まれなければ、
交点長が二以下なのでその逆像は長さ四以下である（各局所環の長さによるフィルトレーションを引き戻せば上界を得る）。
含まれるならGrassmann多様体内の直線は同一平面内の一点を通る直線鉛筆である。
その共通点を $m_0$ とすると $m_0(z)$ は $V(pq)$ 上で消え、非零三次式の次数に反する。
どちらの議論も非被約因子を含む。

ゆえに反対階数九以下で正則なら系\ref{cor:all-defective-one-sym-finite}より両対称である。
両対称表示の有限グラフは転置を交互に繰り返し、残る二位置 $x,y$ の標識付き配置も射影同値となる。
これは非恒等な二十次の根に反する。反対階数十では定理\ref{thm:general-quadric-cone-exclusion}の
固定六次元商の交代形式が全ての $D$ で階数六方向を与え、局所自由性に反する。これで全階数を尽くす。
\end{proof}
\path{computations/verify_cone_fold_plucker.py} はHankel条件の階数と、標本のPlücker五空間・反対切断を再構成する。
